%% Transcription LaTeX du document « CS Finale bis »
%% On fast optimal filtering in Gaussian switching systems — version corrigée et augmentée
\documentclass[12pt, a4paper]{article}

\usepackage[utf8]{inputenc}
\usepackage[T1]{fontenc}
\usepackage[french]{babel}
\usepackage{amsmath, amssymb}
\usepackage{bm}          % \bm{} pour le gras mathématique (section 4)
\usepackage{tikz}
\usetikzlibrary{arrows.meta, positioning}
\usepackage{geometry}
\geometry{top=2.5cm, bottom=2.5cm, left=2.5cm, right=2.5cm}
\usepackage{microtype}

\input{macros}   % \Ek, \pk, \Covk, \Vark, \VarInvk, \Nc, \tp, \inv, \rnn, …

%% Espérance sous le modèle M_j : \EMj{j}{expression}
\newcommand{\EMj}[2]{E_{\mathrm{M}_{#1}}\!\left[#2\right]}

\setlength{\parindent}{0pt}
\setlength{\parskip}{5pt plus 1pt}

\begin{document}

\noindent\hfill\textit{15 avril 26}

\begin{center}
  {\large\bfseries On fast optimal filtering in Gaussian switching systems.}\\[6pt]
  Stéphane Derrode \quad Clément Fernandes \quad Frédéric Lehmann \quad Wojciech Pieczynski
\end{center}

\medskip

Pour $N \geq 3$, soient $X_{1:N} = (X_1, \ldots, X_N)$, $Y_{1:N} = (Y_1, \ldots, Y_N)$, et
$R_{1:N} = (R_1, \ldots, R_N)$, trois processus stochastiques.
$X_{1:N}$ et $Y_{1:N}$ sont à valeurs, respectivement, dans $\mathbb{R}^q$, $\mathbb{R}^s$, et
$R_{1:N}$ est à valeurs dans $\Omega = \{1, \ldots, K\}$.
Le problème est de calculer $\Ek{X_n}{\yN}$, pour tout $n = 1, \ldots, N$.

On sait faire du filtrage et du lissage si
\begin{align}
  \text{(H1)} \quad & (\XN, \RN, \YN)
    \text{ est de Markov et homogène;} \tag{1} \\
  \text{(H2)} \quad & \pk{x_{1:N}, y_{1:N}}{r_{1:N}}
    \text{ est gaussien ;} \tag{2} \\
  \text{(H3)} \quad & \pk{r_{n+1}}{x_n, r_n, y_n} = \pk{r_{n+1}}{r_n}
    \ \text{et}\ \pk{x_{n+1}, y_{n+1}}{r_n, r_{n+1}} = \pk{x_{n+1}, y_{n+1}}{r_{n+1}}
    \tag{3} \\
  \text{(H4)} \quad & \pk{r_{n+1}, y_{n+1}}{x_n, r_n, y_n} = \pk{r_{n+1}, y_{n+1}}{r_n, y_n}
    \tag{4}
\end{align}

Notons qu'en posant $Z_n = [X_n\tp, Y_n\tp]\tp$, la chaîne $(\RN, Z_{1:N})$ est la très classique
HMC à bruit corrélé.

Rappelons que (H4) implique markovianité de $(\RN, \YN)$, qui implique la markovianité de
$\pk{r_{1:N}}{y_{1:N}}$, à l'origine de la possibilité de filtrage optimal rapide..

En fait chacune de ces hypothèses peut probablement être affaiblie, ce qui fait autant d'axes
de recherche possibles. Ici on s'intéresse à (H4). (H1)--(H4) implique que les transitions
sont de la forme
\begin{equation}
  \begin{bmatrix} X_{n+1} \\ Y_{n+1} \end{bmatrix}
  = \begin{bmatrix} A_{r_{n+1}} & B_{r_{n+1}} \\ \mathbf{0} & D_{r_{n+1}} \end{bmatrix}
    \begin{bmatrix} X_n \\ Y_n \end{bmatrix}
  + \begin{bmatrix} U_{n+1} \\ V_{n+1} \end{bmatrix}, \tag{5}
\end{equation}
et l'existence de ce zéro pose problème dans les filtrages avec les modèles à sauts.

Il s'agit donc de trouver une hypothèse différente de (H4), exploitable, qui supprimerait le
zéro en question en gardant la possibilité de filtrage et lissage exacts.

Considérons (H5) (équivalente à (H) de la proposition générale, qui est
$\bm{p(r_n, y_n \mid x_{n+1}, r_{n+1}, y_{n+1}) = p(r_n, y_n|r_{n+1}, y_{n+1})}$) suivante :
\begin{equation}
  \text{(H5)} \quad
  \pk{x_{n+1}}{r_{n+1}, y_{n+1}, r_n, y_n} = \pk{x_{n+1}}{r_{n+1}, y_{n+1}},
  \tag{6}
\end{equation}
équivalent à
\begin{equation}
  \pk{r_n, y_n}{r_{n+1}, y_{n+1}, x_{n+1}} = \pk{r_n, y_n}{r_{n+1}, y_{n+1}}
  \tag{6bis}
\end{equation}

Avec (H1)--(H3) et (H5) le modèle vérifie
\begin{equation}
  \underbrace{\begin{bmatrix} X_{n+1} \\ Y_{n+1} \end{bmatrix}}_{Z_{n+1}}
  = \underbrace{\begin{bmatrix} A_{r_{n+1}} & B_{r_{n+1}} \\
                                C_{r_{n+1}} & D_{r_{n+1}} \end{bmatrix}}_{F(r_{n+1})}
    \underbrace{\begin{bmatrix} X_n \\ Y_n \end{bmatrix}}_{Z_n}
  + {\color{magenta}\underbrace{\begin{bmatrix} b_{X,r_{n+1}} \\ b_{Y,r_{n+1}} \end{bmatrix}}_{b_{r_{n+1}}}}
  + \underbrace{\begin{bmatrix} U_{n+1} \\ V_{n+1} \end{bmatrix}}_{W_{n+1}},
  \tag{7}
\end{equation}
ou encore
\begin{equation}
  Z_{n+1} = F_{(r_{n+1})}\,Z_n + {\color{magenta}b_{r_{n+1}}} + W_{n+1}, \tag{7bis}
\end{equation}
avec
\begin{equation}
  \Vark{W_{n+1}}{r_{n+1}}
  = \begin{bmatrix}
      \Sigma_{U,r_{n+1}} & \Delta_{r_{n+1}} \\
      \Delta_{r_{n+1}}\tp & \Sigma_{V,r_{n+1}}
    \end{bmatrix} \tag{7ter}
\end{equation}
{\color{magenta}et $b_{r_{n+1}} = \bigl[b_{X,r_{n+1}}\tp,\, b_{Y,r_{n+1}}\tp\bigr]\tp \in \mathbb{R}^{q+s}$
le vecteur de biais de dérive dépendant du régime ($b_{r_{n+1}} = 0$ redonne le modèle sans biais).}

\noindent\textbf{Rappel} Dans le cas de mélange gaussien multidimensionnel
$\pu{x, k, y} = \pu{k}\,\pk{x, y}{k}$ on a
$\Ek{X}{y} = \sum_k \pk{k}{y}\,\Ek{X}{k, y}$ et
$\Ek{X_n X_n\tp}{y} = \sum_k \pk{k}{y}\,\Ek{X_n X_n\tp}{k, y}$.
On peut donc calculer $\Ek{X}{y}$ et $\Ek{X_n X_n\tp}{y}$ si on a les $\pk{k}{y}$,
les $\Ek{X}{k, y}$, les $\Ek{X_n X_n\tp}{k, y}$, et on sait calculer les sommes.
Dans le cas qui nous occupe $X = X_n$, $k = r_{1:n} = (r_1, \ldots, r_n)$,
et $y = \yn = (y_1, \ldots, y_n)$. Le filtrage optimal sera « rapide » si on calcule
$\Ek{X_{n+1}}{\ynp}$ et $\Vark{X_{n+1}}{\ynp}$ à partir de $\Ek{X_n}{\yn}$,
$\Vark{X_n}{\yn}$, et $y_{n+1}$ en un nombre fini, indépendant de $n$, d'opérations.
L'étonnant, dû à la markovianité, est que l'on puisse le calculer pour $n$ grand alors même
que, pour $m$ valeurs possibles de $r_n$, $k = m^n$.

Sachant que
\begin{align}
  \Ek{X_{n+1}}{\ynp}
    &= \sum_{r_{n+1}} \pk{r_{n+1}}{\ynp}\,\Ek{X_{n+1}}{r_{n+1}, \ynp} \tag{8} \\[4pt]
  \Vark{X_{n+1}}{\ynp}
    &= \sum_{r_{n+1}} \pk{r_{n+1}}{\ynp}\,\Vark{X_{n+1}}{r_{n+1}, \ynp} \tag{9}
\end{align}
Il suffit de donner le calcul de $\pk{r_{n+1}}{\ynp}$, $\Ek{X_{n+1}}{r_{n+1}, \ynp}$,
et $\Vark{X_{n+1}}{r_{n+1}, \ynp}$ à partir de $\pk{r_n}{\yn}$, $\Ek{X_n}{r_n, \yn}$,
$\Vark{X_n}{y}$, et $y_{n+1}$ en un nombre fini, indépendant de $n$, d'opérations.

\medskip
\noindent\textbf{Remarque.} (H4) et (H5) impliquent que les lois marginales
$\pu{x_n, r_n, y_n}$ de $(\XN, \RN, \YN)$ sont de la forme
$\pu{x_n, r_n, y_n} = \pu{r_n}\,\pk{x_n, y_n}{r_n}$, avec $\pk{x_n, y_n}{r_n}$
gaussiennes, on dira que $(\XN, \RN, \YN)$ est « stationnaire en forme des marginales ».

\bigskip
\noindent\textbf{1. Calcul de $\bm{p(r_{n+1}, y_{n+1}|r_n, y_n)}$ et
  $\bm{p(r_{n+1}|y_{1:n+1})}$}

\medskip
On sait, de par la proposition générale, que $(\RN, \YN)$ est markovien.
Par ailleurs, (H2) et (H3) impliquent que
\begin{equation}
  \pu{x_n, r_n, y_n, x_{n+1}, r_{n+1}, y_{n+1}}
  = \pu{r_n, r_{n+1}}\,\pk{x_n, y_n, x_{n+1}, y_{n+1}}{r_n, r_{n+1}}
  \tag{10}
\end{equation}
avec $\pk{x_n, y_n, x_{n+1}, y_{n+1}}{r_n, r_{n+1}}$ gaussienne. Il en résulte que
\begin{equation}
  \pu{r_n, y_n, r_{n+1}, y_{n+1}}
  = \pu{r_n, r_{n+1}}\,\pk{y_n, y_{n+1}}{r_n, r_{n+1}},
  \tag{11}
\end{equation}
avec $\pk{x_n, y_n, x_{n+1}, y_{n+1}}{r_n, r_{n+1}}$ gaussienne. Il s'ensuit :
\begin{equation}
  \pk{r_{n+1}, y_{n+1}}{r_n, y_n}
  = \pk{r_{n+1}}{r_n}\,\pk{y_{n+1}}{y_n, r_n, r_{n+1}},
  \tag{12}
\end{equation}
avec $\pk{y_{n+1}}{y_n, r_n, r_{n+1}} \sim \Nc(M_{\rnn}\,y_n,\,\Gamma_{\rnn})$
\hfill(13)
\begin{align}
  M_{\rnn} &= \Covk{Y_{n+1}}{Y_n}{\rnn}\,\VarInvk{Y_n}{r_n}
    \tag{14} \\[4pt]
  \Gamma_{\rnn} &= \Vark{Y_{n+1}}{r_{n+1}}
    - \Covk{Y_{n+1}}{Y_n}{\rnn}\,\VarInvk{Y_n}{r_n}\,\Covk{Y_n}{Y_{n+1}}{\rnn}
    \tag{15}
\end{align}

Les matrices de covariances $\Vark{Y_n}{r_n}$, $\Covk{Y_{n+1}}{Y_n}{\rnn}$, et
$\Vark{Y_{n+1}}{r_{n+1}}$ dans (13)--(14) sont extraites des matrices
$\Covk{Z_{n+1}}{Z_n}{\rnn}$ et $\Vark{Z_{n+1}}{r_{n+1}}$
--- on pose $Z_n = [X_n\tp, Y_n\tp]\tp$, voir (7bis) --- calculées récursivement par
\begin{equation}
  \Covk{Z_{n+1}}{Z_n}{\rnn} = F_{r_{n+1}}\,\Vark{Z_n}{r_n}; \tag{16}
\end{equation}
\begin{equation}
  \Vark{Z_{n+1}}{r_{n+1}}
  = \sum_{r_n} \pk{r_n}{r_{n+1}}\,F_{r_{n+1}}\,\Vark{Z_n}{r_n}\,F_{r_{n+1}}\tp
    + \Vark{W_{n+1}}{r_{n+1}}, \tag{17}
\end{equation}
avec
\begin{equation}
  \pk{r_n}{r_{n+1}} = \frac{\pu{r_n}\,\pk{r_{n+1}}{r_n}}%
    {\displaystyle\sum_{r_n} \pu{r_n}\,\pk{r_{n+1}}{r_n}}
  \tag{17bis}
\end{equation}
et les marginales $\pu{r_n}$ calculées classiquement.

{\color{green!60!black}
\smallskip
\noindent\textbf{Option B — moyenne non nulle.}
Si $\mu_n(r) := \Ek{Z_n}{r_n=r} \neq 0$, les seconds moments $\Vark{Z_n}{r_n}$
seront notés $P_n(r_n)$ et les covariances \emph{centrées} $\Sigma_n(r_n) := P_n(r_n) - \mu_n(r_n)\,\mu_n(r_n)\tp$
remplacent les seconds moments dans toutes les formules de conditionnement.
Les moyennes marginales se propagent en parallèle des seconds moments (17) :
\begin{equation}
  \mu_{n+1}(r_{n+1}) \;=\; F_{r_{n+1}}\!\sum_{r_n} \pk{r_n}{r_{n+1}}\,\mu_n(r_n)
  {\color{magenta}\;+\; b_{r_{n+1}}}.
  \tag{17ter}
\end{equation}
En particulier, l'équation (13) se corrige en
\begin{equation}
  \pk{y_{n+1}}{y_n,\,\rnn}
  \;\sim\; \Nc\!\left(
    \mu_{Y,n+1}(\rnn)
    + \widetilde{M}_{\rnn}\,\bigl(y_n - \mu_{Y,n}(r_n)\bigr),\;
    \widetilde{\Gamma}_{\rnn}
  \right),
  \tag{13'}
\end{equation}
où $\mu_{Y,n+1}(\rnn) := \bigl[C_{r_{n+1}},\,D_{r_{n+1}}\bigr]\,\mu_n(r_n)
{\color{magenta}\;+\; b_{Y,r_{n+1}}}$
est la moyenne de $Y_{n+1}$ sachant $\rnn$, et $\widetilde{M}_{\rnn}$, $\widetilde{\Gamma}_{\rnn}$
sont définis par (14)--(15) en substituant les blocs de $\Sigma_n(\cdot)$ aux blocs de $P_n(\cdot)$.
}

Cela termine le calcul de $\pk{r_{n+1}, y_{n+1}}{r_n, y_n}$ ; $\pk{r_{n+1}}{\ynp}$ est
donné classiquement par :
\begin{equation}
  \pk{r_{n+1}}{\ynp}
  = \frac{\displaystyle\sum_{r_n} \pu{r_n, \yn}\,\pk{r_{n+1}, y_{n+1}}{r_n, y_n}}
         {\displaystyle\sum_{r_n, r_{n+1}} \pu{r_n, \yn}\,\pk{r_{n+1}, y_{n+1}}{r_n, y_n}}
  \tag{18}
\end{equation}
En effet :
\begin{align*}
  \pk{r_{n+1}}{\ynp}
  &= \sum_{r_n} \pk{r_n, r_{n+1}}{\ynp} \\
  &= \frac{\displaystyle\sum_{r_n} \pu{r_n, r_{n+1}, \ynp}}
          {\displaystyle\sum_{r_n, r_{n+1}} \pu{r_n, r_{n+1}, \ynp}}
   = \frac{\displaystyle\sum_{r_n} \pu{r_n, \yn}\,\pk{r_{n+1}, y_{n+1}}{r_n, y_n}}
          {\displaystyle\sum_{r_n, r_{n+1}} \pu{r_n, \yn}\,\pk{r_{n+1}, y_{n+1}}{r_n, y_n}}
\end{align*}

\bigskip
\noindent\textbf{2. Calcul de $\bm{E[X_{n+1}|r_{n+1}, y_{1:n+1}]}$ et
  $\bm{E[X_{n+1}X_{n+1}^T|r_{n+1}, y_{1:n+1}]}$}

\medskip
Le calcul de $\Ek{X_{n+1}}{r_{n+1}, \ynp}$ est immédiat à partir des identités suivantes,
justifiée (en bleu) juste après le point~3 :
\begin{align}
  \Ek{X_{n+1}}{r_{n+1}, \ynp} &= \Ek{X_{n+1}}{r_{n+1}, y_{n+1}} \tag{19} \\[4pt]
  \Vark{X_{n+1}}{r_{n+1}, \ynp} &= \Vark{X_{n+1}}{r_{n+1}, y_{n+1}} \tag{20}
\end{align}
En effet, $\Ek{X_{n+1}}{r_{n+1}, y_{n+1}}$ et $\Vark{X_{n+1}}{r_{n+1}, \ynp}$ sont
calculées simplement à partir de $\Vark{Z_{n+1}}{r_{n+1}}$, donnée séquentiellement par
(16) et (17) :
\[
  \pk{x_{n+1}}{y_{n+1}, r_{n+1}}
  \sim \Nc\!\left(M_{r_{n+1}}\,y_{n+1},\,\Gamma_{r_{n+1}}\right),
\]
avec
\begin{align}
  M_{r_{n+1}} &= \Covk{X_{n+1}}{Y_{n+1}}{r_{n+1}}\,\VarInvk{Y_{n+1}}{r_{n+1}}
    \tag{21} \\[4pt]
  \Gamma_{\rnn} &= \Vark{X_{n+1}}{r_{n+1}}
    - \Covk{X_{n+1}}{Y_{n+1}}{r_{n+1}}\,\VarInvk{Y_{n+1}}{r_{n+1}}\,\Covk{Y_{n+1}}{X_{n+1}}{r_{n+1}}
    \tag{22}
\end{align}

{\color{green!60!black}
\smallskip
\noindent\textbf{Option B.}
La distribution $\pk{x_{n+1}}{r_{n+1}, y_{n+1}} \sim \Nc(M_{r_{n+1}}\,y_{n+1},\,\Gamma_{r_{n+1}})$
ci-dessus suppose $\mu_{n+1}(r_{n+1}) = 0$.
Dans le cas général, le gain de Kalman (21) utilise les blocs de la covariance centrée
$\Sigma_{n+1}(r_{n+1}) = P_{n+1}(r_{n+1}) - \mu_{n+1}(r_{n+1})\,\mu_{n+1}(r_{n+1})\tp$,
et la moyenne conditionnelle se corrige en
\begin{equation}
  \Ek{X_{n+1}}{r_{n+1},\,y_{n+1}}
  \;=\; \mu_{X,n+1}(r_{n+1})
        + K_{r_{n+1}}\,\bigl(y_{n+1} - \mu_{Y,n+1}(r_{n+1})\bigr),
  \tag{21'}
\end{equation}
où $K_{r_{n+1}} = \Sigma_{XY,n+1}(r_{n+1})\,\Sigma_{YY,n+1}(r_{n+1})\inv$,
et $\mu_{X,n+1}(r_{n+1})$, $\mu_{Y,n+1}(r_{n+1})$ sont les blocs $X$ et $Y$
de $\mu_{n+1}(r_{n+1})$ calculé par (17ter).
La formule (22) pour $\Gamma_{r_{n+1}}$ reste inchangée en substituant $\Sigma_{n+1}$ à $P_{n+1}$.
}

\bigskip
\noindent\textbf{3. Calcul de $\bm{E[X_{n+1}|y_{1:n+1}]}$ et
  $\bm{E[X_{n+1}X_{n+1}^T|y_{1:n+1}]}$ par (8)--(9) et (21)--(22)}

\medskip
{\color{blue}
Pour justifier (19), (20) on reprend la démonstration de la proposition~1.
En posant $w_n = (r_n, y_n)$ pour tout $n$, on a, après l'inversion du temps
(on prend $N = n + 1$) :
\begin{align*}
  \pu{x_{1:n+1}, w_{1:n+1}} = {}
  &\bigl[\pk{w_1}{w_2}\,\pk{w_2}{w_3}\cdots\pk{w_{n-1}}{w_n}\,\pk{w_n}{w_{n+1}}\,\pu{w_{n+1}}\bigr] \\
  &\bigl[\pk{x_1}{w_1, w_2, x_2}\,\pk{x_2}{w_2, w_3, x_3}\cdots \\
  &\quad (x_{n-1}|w_{n-1}, w_n, x_n)\,\pk{x_n}{w_n, w_{n+1}, x_{n+1}}\,\pk{x_{n+1}}{w_{n+1}}\bigr].
\end{align*}
En intégrant par rapport à $x_1, x_2, \ldots, x_n$, cela donne
\[
  \pu{x_{n+1}, w_{1:n+1}}
  = \pk{w_1}{w_2}\,\pk{w_2}{w_3}\cdots\pk{w_{n-1}}{w_n}\,\pk{w_n}{w_{n+1}}\,\pu{w_{n+1}}\,
    \pk{x_{n+1}}{w_{n+1}},
\]
qui implique $\pk{x_{n+1}}{w_{1:n+1}} = \pk{x_{n+1}}{w_{n+1}}$, d'où (19) et (20).
}% fin {\color{blue}}

\bigskip
\noindent\textbf{En résumé, le filtrage est le suivant :}

\medskip
Soient $\pk{r_n}{\yn}$, $\Ek{Z_n}{r_n, \yn}$,
$\Vark{Z_n}{\yn} = \Vark{Z_n}{y_n}$ données
(qui donnent donc également $\Ek{X_n}{r_n, y_n}$ et $\Vark{X_n}{\yn}$).
{\color{green!60!black}En option B, l'état inclut aussi $\mu_n(r_n) = \Ek{Z_n}{r_n}$ pour chaque $r_n$.}

\begin{itemize}
  \item[(i)] calcul de $\Covk{Z_{n+1}}{Z_n}{\rnn}$ et $\Vark{Z_{n+1}}{r_{n+1}}$
    avec (16)--(17), ensuite $\pk{r_{n+1}, y_{n+1}}{r_n, y_n}$ avec (12)--(15), et
    $\pk{r_{n+1}}{\ynp}$ avec (18) ;
    {\color{green!60!black}[option B : calculer aussi $\mu_{n+1}(r_{n+1})$ par (17ter) {\color{magenta}(avec biais $b_{r_{n+1}}$)} et utiliser (13') à la place de (13)]}
  \item[(ii)] calcul de $\Ek{X_{n+1}}{r_{n+1}, \ynp}$
    (qui est égal à $\Ek{X_{n+1}}{r_{n+1}, y_{n+1}}$) et
    $\Vark{X_{n+1}}{r_{n+1}, \ynp}$
    (qui est égal à $\Vark{X_{n+1}}{r_{n+1}, y_{n+1}}$) avec (19)--(22) ;
    {\color{green!60!black}[option B : utiliser (21') à la place de $M_{r_{n+1}} y_{n+1}$]}
  \item[(iii)] calcul de $\Ek{X_{n+1}}{\ynp}$ et $\Vark{X_{n+1}}{\ynp}$ par (8)--(9).
\end{itemize}

\bigskip

% -----------------------------------------------------------------------
% Figure 1 — Calcul récursif du filtre optimal
% -----------------------------------------------------------------------
\begin{figure}[h!]
\centering
\begin{tikzpicture}[
  box/.style  = {draw, rounded corners=3pt, minimum width=6.0cm,
                 minimum height=1.4cm, align=center, font=\small},
  lbox/.style = {draw, rounded corners=3pt, minimum width=5.2cm,
                 minimum height=1.4cm, align=center, font=\small},
  arr/.style  = {-{Stealth[length=7pt, width=5pt]}, thick},
  node distance = 1.0cm and 2.0cm
]

%% Boîte gauche (unique)
\node[lbox] (L) {%
  $\Ek{Z_n}{r_n,\,\yn}$,\quad $\Vark{Z_n}{\yn}$};

%% Colonne droite (3 boîtes)
\node[box, right=of L] (R1) {%
  $\Ek{Z_{n+1}}{r_{n+1}}$,\quad $\Covk{Z_{n+1}}{Z_n}{\rnn}$,\\[2pt]
  $\Vark{Z_{n+1}}{r_{n+1}}$};

\node[box, below=of R1] (R2) {%
  $\pk{r_{n+1}, y_{n+1}}{r_n, y_n}$,\quad $\pk{r_{n+1}}{\ynp}$};

\node[box, below=of R2] (R3) {%
  $\Ek{X_{n+1}}{\ynp}$,\quad $\Vark{X_{n+1}}{\ynp}$};

%% Flèche horizontale
\draw[arr] (L) -- (R1);

%% Flèches verticales (colonne droite)
\draw[arr] (R1) -- (R2);
\draw[arr] (R2) -- (R3);

\end{tikzpicture}
\caption{Calcul récursif du filtre optimal et de son erreur quadratique}
\end{figure}

Finalement, l'important est les récursions (16) et (17), qui permettent de calculer les lois
des marginales $(X_n, R_n, Y_n)$ et la loi de $(\RN, \YN)$.
Ensuite (19) et (20) permettent de conclure.

% -----------------------------------------------------------------------
% Section 0 — Initialisation du filtre (n = 1)
% -----------------------------------------------------------------------
\bigskip
\noindent\textbf{0.\ Initialisation du filtre ($n = 1$)}

\medskip
La récursion (i)--(iii) démarre à $n = 1$ après réception de $y_1$.
Le modèle fournit, pour chaque régime $r \in \{1, \ldots, K\}$ :
\begin{itemize}
  \item[$-$] $\pi_0(r) = p(r_1 = r)$ : distribution stationnaire de $P$ ;
  \item[$-$] $\mu_{z_0}(r) = \Ek{Z_1}{r_1=r} \in \mathbb{R}^{q+s}$ ;
  \item[$-$] $\Sigma_{z_0}(r) = \mathrm{Var}(Z_1 \mid r_1=r) \in \mathbb{S}_{++}^{q+s}$.
\end{itemize}
On note les blocs de taille $q$ et $s$ :
\[
  \mu_{z_0}(r) = \begin{bmatrix}\mu_{X,0}(r)\\\mu_{Y,0}(r)\end{bmatrix}, \qquad
  \Sigma_{z_0}(r) = \begin{bmatrix}\Sigma_{XX,0}(r) & \Sigma_{XY,0}(r) \\
    \Sigma_{XY,0}(r)\tp & \Sigma_{YY,0}(r)\end{bmatrix}.
\]
Le second moment initial --- nécessaire pour les équations (16)--(17) --- est :
\begin{equation}
  P_{z_0}(r) \;:=\; \Ek{Z_1 Z_1\tp}{r_1=r}
  \;=\; \Sigma_{z_0}(r) + \mu_{z_0}(r)\,\mu_{z_0}(r)\tp
  \;=\; \begin{bmatrix}P_{XX,0}(r) & P_{XY,0}(r)\\P_{XY,0}(r)\tp & P_{YY,0}(r)\end{bmatrix}.
  \tag{I.1}
\end{equation}

\noindent\textbf{Traitement de la première observation $y_1$.}
La loi marginale de $Y_1$ sachant $r_1 = r$ est
$\pk{y_1}{r_1=r} = \Nc\!\left(y_1;\,\mu_{Y,0}(r),\,\Sigma_{YY,0}(r)\right)$, d'où
\begin{equation}
  \pk{r_1=r}{y_1} = \frac{\pi_0(r)\;\Nc\!\left(y_1;\,\mu_{Y,0}(r),\,\Sigma_{YY,0}(r)\right)}%
    {\displaystyle\sum_{r'=1}^{K}\pi_0(r')\;\Nc\!\left(y_1;\,\mu_{Y,0}(r'),\,\Sigma_{YY,0}(r')\right)}.
  \tag{I.2}
\end{equation}
Le gain de Kalman initial et la moyenne conditionnelle de $Z_1$ sachant $(r_1=r,\,y_1)$ sont :
\begin{equation}
  M_r = \Sigma_{XY,0}(r)\,\Sigma_{YY,0}(r)\inv, \qquad
  \Ek{Z_1}{r_1=r,\,y_1}
  = \begin{bmatrix}
      \mu_{X,0}(r) + M_r\bigl(y_1 - \mu_{Y,0}(r)\bigr)\\[3pt] y_1
    \end{bmatrix}.
  \tag{I.3}
\end{equation}
Les grandeurs $\pk{r_1}{y_1}$, $\Ek{Z_1}{r_1,\,y_1}$ et $P_{z_0}(r)$ constituent l'état initial
de la récursion (i)--(iii).
{\color{green!60!black}En option B, $\mu_1(r) := \mu_{z_0}(r)$ est la valeur initiale de la récursion (17ter) ;
la formule (I.3) est déjà la formule générale avec correction de moyenne.}

% -----------------------------------------------------------------------
% Section 4 — Simulations (en rouge : rédaction préliminaire)
% -----------------------------------------------------------------------
\bigskip
{\color{red}
\noindent\textbf{4. Simulations}

\medskip
Attention, c'est plus une idée générale qu'une rédaction rigoureuse.

Considérons trois modèles triplets $T_{1:\mathrm{N}} = (\mathbf{X_{1:N}}, \mathbf{R_{1:N}},
\mathbf{Y_{1:N}})$, notés $\mathrm{M}_1$, $\mathrm{M}_2$, $\mathrm{M}_3$, vérifiant
(H1)--(H3). Leurs lois sont définies par les lois $\pk{r_1, r_2}{}$%
\footnote{On utilise ici la notation par paires $(r_1,r_2)$, $(z_1,z_2)$ pour la présentation
  à deux pas de temps.}
et $\pk{z_1, z_2}{r_1, r_2}$. En effet, on peut, à partir de ces lois données, déterminer la
présentation (7)--(7ter).

On prend les $\pu{r_1, r_2}$, $\pk{z_1}{r_1}$, et $\pk{z_2}{r_2}$ communes aux trois modèles,
ainsi que les covariances $\Covk{X_2}{X_1}{r_1, r_2}$, et $\Covk{Y_2}{Y_1}{r_1, r_2}$.
La différence entre $\mathrm{M}_1$ et $\mathrm{M}_2$ (qui sera de type « $C=0$ ») est au
niveau de $\Covk{Y_2}{X_1}{r_2, r_1}$, et la différence entre $\mathrm{M}_1$ et $\mathrm{M}_3$
(qui sera le nouveau modèle proposé) est au niveau de $\Covk{Y_1}{X_2}{r_1, r_2}$.
L'idée est donc de considérer $\mathrm{M}_2$ et $\mathrm{M}_3$ comme des
« approximations » de $\mathrm{M}_1$, et comparer leurs qualités en tant que telles.

La loi de $\mathrm{M}_1$ est donc donnée par $\pu{r_1, r_2}$, $\pk{z_1}{r_1}$, et
$\pk{z_2}{r_2}$, et
\begin{equation}
  \EMj{1}{Z_2 Z_1\tp \mid r_1, r_2}
  = \begin{bmatrix}
      \Covk{X_2}{X_1}{r_2, r_1} & \Covk{X_2}{Y_1}{r_2, r_1} \\[4pt]
      \Covk{Y_2}{X_1}{r_2, r_1} & \Covk{Y_2}{Y_1}{r_2, r_1}
    \end{bmatrix}
  \tag{4.1}
\end{equation}

La loi de $\mathrm{M}_2$ est obtenue à partir de (4.1) en remplaçant
$\Covk{Y_2}{X_1}{r_2, r_1}$ par
$\Covk{Y_2}{Y_1}{r_2, r_1}\,\mathbf{E}^{-1}[\mathbf{Y_1 Y_1\tp}\,|\,r_1]\,\mathbf{E}[\mathbf{Y_1 X_1\tp}\,|\,\mathbf{r_1}]$ :
\begin{equation}
  \EMj{2}{Z_2 Z_1\tp \mid r_1, r_2}
  = \begin{bmatrix}
      \Covk{X_2}{X_1}{r_2, r_1} & \Covk{X_2}{Y_1}{r_2, r_1} \\[4pt]
      \Covk{Y_2}{Y_1}{r_2, r_1}\,\mathbf{E}^{-1}[\mathbf{Y_1 Y_1\tp}\,|\,r_1]\,\mathbf{E}[\mathbf{Y_1 X_1\tp}\,|\,\mathbf{r_1}]
        & \Covk{Y_2}{Y_1}{r_2, r_1}
    \end{bmatrix}
  \tag{4.2}
\end{equation}

Et la loi de $\mathrm{M}_3$, obtenue à partir de $\mathrm{M}_1$ en remplaçant l'espérance
$\Covk{Y_1}{X_2}{r_1, r_2}$ par
$\Covk{Y_1}{Y_2}{r_1, r_2}\,\mathbf{E}^{-1}[\mathbf{Y_2 Y_2\tp}\,|\,r_2]\,\mathbf{E}[\mathbf{Y_2 X_2\tp}\,|\,\mathbf{r_2}]$,
ce qui revient à remplacer dans (4.1) $\Covk{X_2}{Y_1}{r_2, r_1}$ par
$\Covk{X_2}{Y_2}{r_2}\,E\inv[\Covk{Y_2}{Y_2}{r_2}]\,\Covk{Y_2}{Y_1}{r_2, r_1}$ :
\begin{equation}
  \EMj{3}{Z_2 Z_1\tp \mid r_1, r_2}
  = \begin{bmatrix}
      \Covk{X_2}{X_1}{r_2, r_1}
        & \Covk{X_2}{Y_2}{r_2}\,\VarInvk{Y_2}{r_2}\,\Covk{Y_2}{Y_1}{r_2, r_1} \\[4pt]
      \Covk{Y_2}{X_1}{r_2, r_1} & \Covk{Y_2}{Y_1}{r_2, r_1}
    \end{bmatrix}
  \tag{4.3}
\end{equation}

Ainsi, en considérant le modèle « $C=0$ » et le nouveau proposé comme deux approximations
de $\mathrm{M}_1$, on pourrait évaluer laquelle est meilleure (en veillant que les deux
approximations soient du « même ordre »).

Notons que (4.3) donne $F_{(r_{n+1})}$ de $\mathrm{M}_3$ :
\begin{equation*}
  \mathbf{F}_{\mathrm{M}_3,\, r_{n+1}}
  = \EMj{3}{Z_2 Z_1\tp \mid r_1, r_2}\,\VarInvk{Y_1}{r_1}
\end{equation*}
On devrait donc pouvoir simuler selon $\mathrm{M}_3$ et filtrer selon $\mathrm{M}_3$, et
réciproquement. Ce type d'étude pourrait servir --- par ailleurs --- aux études du lissage.

\bigskip
\noindent\textbf{Contrainte H5 sur les paramètres — déduction de $B$.}

\smallskip
La condition (H5) impose une relation algébrique entre les 7 matrices du modèle
($A$, $B$, $C$, $D$, $\Sigma_U$, $\Sigma_V$, $\Delta$, toutes indexées par le régime courant).
En omettant provisoirement le subscript $r_{n+1}$, elle s'écrit :
\begin{equation}
  \Delta^T A + \Sigma_V B^T
  = (\Delta^T C^T + \Sigma_V D^T)
    \bigl[(C\Sigma_U + D\Delta^T)C^T + (C\Delta + D\Sigma_V)D^T + \Sigma_V\bigr]^{-1}
    \bigl[(C\Sigma_U + D\Delta^T)A^T + (C\Delta + D\Sigma_V)B^T + \Delta^T\bigr].
  \tag{4.4}
\end{equation}
En réintroduisant le subscript $r_{n+1}$ sur chacune des 7 matrices, on obtient la forme
complète :
\begin{align}
  \Delta_{r_{n+1}}^T A_{r_{n+1}} + \Sigma_{V,r_{n+1}} B_{r_{n+1}}^T
  &= \bigl(\Delta_{r_{n+1}}^T C_{r_{n+1}}^T + \Sigma_{V,r_{n+1}} D_{r_{n+1}}^T\bigr)
  \notag\\
  &\quad\times
    \bigl[
      (C_{r_{n+1}}\Sigma_{U,r_{n+1}} + D_{r_{n+1}}\Delta_{r_{n+1}}^T)\,C_{r_{n+1}}^T
     +(C_{r_{n+1}}\Delta_{r_{n+1}} + D_{r_{n+1}}\Sigma_{V,r_{n+1}})\,D_{r_{n+1}}^T
     + \Sigma_{V,r_{n+1}}
    \bigr]^{-1}
  \notag\\
  &\quad\times
    \bigl[
      (C_{r_{n+1}}\Sigma_{U,r_{n+1}} + D_{r_{n+1}}\Delta_{r_{n+1}}^T)\,A_{r_{n+1}}^T
     +(C_{r_{n+1}}\Delta_{r_{n+1}} + D_{r_{n+1}}\Sigma_{V,r_{n+1}})\,B_{r_{n+1}}^T
     + \Delta_{r_{n+1}}^T
    \bigr].
  \tag{4.5}
\end{align}

\medskip
\noindent\textbf{Déduction de $B$ — système linéaire.}
L'équation (4.4) est \emph{linéaire} en $B^T$ : elle apparaît des deux membres, mais avec des
coefficients matriciels connus dès que les 6 autres matrices sont fixées. On introduit les
notations compactes (subscript $r$ omis) :
\begin{equation}
  P = \Delta^T C^T + \Sigma_V D^T, \qquad
  Q = C\Sigma_U + D\Delta^T,       \qquad
  R = C\Delta + D\Sigma_V,         \qquad
  M = QC^T + RD^T + \Sigma_V.
  \tag{4.6}
\end{equation}
\textit{Remarque.} $R = P^T$ ($\Sigma_V$ symétrique) et $M = M^T \succ 0$ dès que
$\Sigma_U, \Sigma_V \succ 0$.

\noindent L'équation (4.4) se réécrit $\Delta^T A + \Sigma_V B^T = PM^{-1}(QA^T + RB^T +
\Delta^T)$. En regroupant les termes en $B^T$ dans le membre gauche :
\begin{equation}
  \underbrace{\bigl(\Sigma_V - PM^{-1}R\bigr)}_{\displaystyle L}\; B^T
  \;=\; PM^{-1}(QA^T + \Delta^T) - \Delta^T A.
  \tag{4.7}
\end{equation}
La matrice $L = \Sigma_V - PM^{-1}R = \Sigma_V - PM^{-1}P^T$ est le
\emph{complément de Schur} de $M$ dans la matrice-bloc symétrique
$\begin{bmatrix} M & P^T \\ P & \Sigma_V \end{bmatrix}$.
Si $M \succ 0$, alors $L \succ 0$ et l'unique solution est :
\begin{equation}
  \boxed{B^T \;=\; L^{-1}\!\bigl(PM^{-1}(QA^T + \Delta^T) - \Delta^T A\bigr).}
  \tag{4.8}
\end{equation}
Ainsi, dès que $A$, $C$, $D$, $\Sigma_U$, $\Sigma_V$, $\Delta$ sont fixées, la matrice $B$ est
\emph{entièrement déterminée} par (4.8) : la contrainte (H5) lève le degré de liberté
correspondant.

\medskip
\noindent\textbf{Déduction de $A$ — système linéaire.}
La même équation (4.7) peut être réorganisée pour exhiber $A^T$ comme inconnue, lorsque
$B$, $C$, $D$, $\Sigma_U$, $\Sigma_V$, $\Delta$ sont fixées. En isolant les deux termes
en $A^T$ (celui dans $QA^T$ à droite et celui dans $\Delta^T A^T$ qui provient de $\Delta^T A$
dans la version transposée de la contrainte) :
\begin{equation}
  P M^{-1}\,Q\,A^T \;-\; \Delta^T A^T
  \;=\; L B^T \;-\; P M^{-1}\,\Delta^T,
  \tag{4.9}
\end{equation}
soit, en factorisant à gauche :
\begin{equation}
  \underbrace{\bigl(P M^{-1} Q - \Delta^T\bigr)}_{\displaystyle G}\; A^T
  \;=\; L\,B^T \;-\; P M^{-1}\,\Delta^T,
  \tag{4.10}
\end{equation}
où $G \in \mathbb{R}^{s \times q}$ et le second membre est dans $\mathbb{R}^{s \times q}$.
La résolution dépend du rapport entre $s$ et $q$ :
\begin{itemize}
  \item \textbf{si $s \geq q$ et $G$ est de rang plein $q$,} la solution unique au sens
    des moindres carrés (qui annule le résidu) est :
    \begin{equation}
      \boxed{A^T \;=\; \bigl(G^T G\bigr)^{-1}\,G^T\,\bigl(L B^T - P M^{-1} \Delta^T\bigr) ;}
      \tag{4.11}
    \end{equation}
  \item \textbf{si $s < q$ ou $G$ est rang-déficient,} $A$ n'est pas uniquement déterminé
    par (H5) : il subsiste $q - \mathrm{rg}(G) \geq 1$ degrés de liberté.
\end{itemize}
En pratique, on résout (4.10) au moyen de \texttt{numpy.linalg.lstsq} et on signale
toute déficience de rang.

\medskip
\noindent\textbf{Déduction de $\Sigma_U$ — système linéaire vectorisé.}
Lorsque $A$, $B$, $C$, $D$, $\Delta$, $\Sigma_V$ sont fixées et que l'on cherche $\Sigma_U$,
le problème devient \emph{non linéaire} dans la forme brute de (4.4) car $\Sigma_U$
apparaît à la fois dans $Q$ ($Q = C\Sigma_U + D\Delta^T$) et dans $M$
($M = QC^T + RD^T + \Sigma_V$ contient le terme $C\Sigma_U C^T$). On
\emph{linéarise} en multipliant (4.7) à gauche par $M$, ce qui élimine $M^{-1}$.

\noindent Posons les blocs « libres de $\Sigma_U$ » :
\begin{equation}
  Q_0 := D\,\Delta^T,                  \qquad
  M_0 := Q_0\,C^T + R\,D^T + \Sigma_V, \qquad
  R   := C\,\Delta + D\,\Sigma_V,
  \tag{4.12}
\end{equation}
de sorte que $Q = Q_0 + C\,\Sigma_U$ et $M = M_0 + C\,\Sigma_U\,C^T$. En multipliant
(4.7) par $M$ à gauche et en regroupant les termes :
\begin{equation}
  \bigl(M_0 + C\,\Sigma_U\,C^T\bigr)\bigl(\Delta^T A + \Sigma_V B^T\bigr)
  \;=\; P\,\bigl((Q_0 + C\,\Sigma_U)\,A^T + R\,B^T + \Delta^T\bigr).
  \tag{4.13}
\end{equation}
En notant
\begin{equation}
  Z := \Sigma_V\,B^T + \Delta^T A,    \qquad
  W := Q_0\,A^T + R\,B^T + \Delta^T,
  \tag{4.14}
\end{equation}
(tous deux de taille $s \times q$ et indépendants de $\Sigma_U$), l'égalité (4.13) se
réécrit :
\begin{equation}
  C\,\Sigma_U\,C^T\,Z \;-\; P\,C\,\Sigma_U\,A^T
  \;=\; \underbrace{P\,W \;-\; M_0\,Z}_{\displaystyle \mathrm{RHS}\,\in\,\mathbb{R}^{s\times q}}.
  \tag{4.15}
\end{equation}

\noindent Pour transformer cette équation matricielle en un système linéaire standard, on
utilise l'identité de \emph{vectorisation}
$\mathrm{vec}(XYZ) = (Z^T \otimes X)\,\mathrm{vec}(Y)$ (où $\otimes$ est le produit de
Kronecker). En posant $E := C^T Z$ ($q \times q$) et $PC := P\,C$ ($s \times q$), il vient
$\mathrm{vec}(C \Sigma_U C^T Z) = (E^T \otimes C)\,\mathrm{vec}(\Sigma_U)$ et
$\mathrm{vec}(P C \Sigma_U A^T) = (A \otimes PC)\,\mathrm{vec}(\Sigma_U)$, d'où le
système linéaire $(sq) \times q^2$ :
\begin{equation}
  \boxed{\bigl[(E^T \otimes C) - (A \otimes PC)\bigr]\,\mathrm{vec}(\Sigma_U)
  \;=\; \mathrm{vec}(\mathrm{RHS}).}
  \tag{4.16}
\end{equation}
La résolution se fait par moindres carrés ; $\Sigma_U$ est ensuite \emph{symétrisée}
($\Sigma_U \leftarrow (\Sigma_U + \Sigma_U^T)/2$) puis on vérifie sa
\emph{définie-positivité} par décomposition de Cholesky.

\medskip
\noindent\textbf{Bilan.} La contrainte (H5) lie les sept matrices $(A, B, C, D, \Sigma_U,
\Delta, \Sigma_V)$ d'un même régime par une seule équation matricielle. Selon la matrice
que l'on considère comme inconnue :
\begin{center}
\begin{tabular}{@{}lll@{}}
  \textbf{Inconnue} & \textbf{Système (équation)} & \textbf{Existence/unicité} \\ \hline
  $B$       & $L\,B^T = \mathrm{rhs}_B$        \quad (4.7)  & unique si $L \succ 0$ \\
  $A$       & $G\,A^T = \mathrm{rhs}_A$        \quad (4.10) & unique si $\mathrm{rg}(G) = q$ \\
  $\Sigma_U$ & Kronecker $(sq \times q^2)$    \quad (4.16) & unique si $\mathrm{rg} = q^2$, puis SPD \\
\end{tabular}
\end{center}

Les trois projections sont implémentées dans \texttt{prg/utils/h5\_constraint.py}
(\texttt{compute\_B\_from\_h5}, \texttt{compute\_A\_from\_h5},
\texttt{compute\_SU\_from\_h5}), avec contrôle de conditionnement et levée d'exception
en cas de système mal posé.

}% fin {\color{red}}

\end{document}
